\documentclass[conference]{IEEEtran}
\IEEEoverridecommandlockouts
\usepackage{cite}
\usepackage{amsmath,amssymb,amsfonts,amsthm}
\usepackage{algorithmic}
\usepackage{graphicx}
\graphicspath{{figures/}{./}} % Fallback to root if user uploads images directly
\usepackage{textcomp}
\usepackage{xcolor}
\usepackage{booktabs}
\usepackage{multirow}
\usepackage{array}
\usepackage{url}
\usepackage[hidelinks]{hyperref}

\newtheorem{theorem}{Theorem}
\newtheorem{remark}{Remark}
\def\BibTeX{{\rm B\kern-.05em{\sc i\kern-.025em b}\kern-.08em
    T\kern-.1667em\lower.7ex\hbox{E}\kern-.125emX}}
\raggedbottom % Prevents huge vertical gaps between paragraphs
% Tighter float spacing
\setlength{\textfloatsep}{10pt plus 2pt minus 4pt}
\setlength{\floatsep}{10pt plus 2pt minus 4pt}
\setlength{\intextsep}{10pt plus 2pt minus 4pt}
\begin{document}

\title{Graph Neural Networks for Learned Matrix Chain Optimization:\\When Structural Reasoning Outperforms Statistical Estimation}

\author{
\IEEEauthorblockN{1\textsuperscript{st} Ch.\ Abhilash}
\IEEEauthorblockA{\textit{Department of CSE} \\
\textit{VNR VJIET}\\
Hyderabad, India \\
abhilash00042@gmail.com}
\and
\IEEEauthorblockN{2\textsuperscript{nd} Dakshatha}
\IEEEauthorblockA{\textit{Department of CSE} \\
\textit{VNR VJIET}\\
Hyderabad, India \\
dklakshatha18@gmail.com}
\and
\IEEEauthorblockN{3\textsuperscript{rd} Shaun Angel}
\IEEEauthorblockA{\textit{Department of CSE} \\
\textit{VNR VJIET}\\
Hyderabad, India \\
dshaunangel@gmail.com}
\and
\IEEEauthorblockN{4\textsuperscript{th} Abhinav}
\IEEEauthorblockA{\textit{Department of CSE} \\
\textit{VNR VJIET}\\
Hyderabad, India \\
abhinavkanna2627@gmail.com}
\and
\IEEEauthorblockN{5\textsuperscript{th} Maneeshwar}
\IEEEauthorblockA{\textit{Department of CSE} \\
\textit{VNR VJIET}\\
Hyderabad, India \\
maneeshwarkarne@gmail.com}
}

\maketitle

\begin{abstract}
The Matrix Chain Multiplication (MCM) problem requires finding the most efficient
sequence for multiplying a chain of matrices. Classical dynamic programming (DP)
solves it exactly in $O(n^3)$ time, which becomes a practical bottleneck in systems
that must evaluate many chains repeatedly, such as deep learning compilers and
database query optimizers. This paper studies whether learned models can
approximate the DP solver with low error, and---a central finding is that---
\emph{how the problem is framed} determines whether learned predictions are even
mathematically meaningful. Two paradigms are compared under identical conditions:
\emph{structural reasoning}, in which a Graph Neural Network (GraphMCMNet) and a
Pointer Network (PointerMCMNet) predict the DP split table and the cost is then
computed exactly from the predicted structure, versus \emph{statistical
estimation}, in which XGBoost and Random Forest ensembles regress the scalar cost
from 145 engineered features. We prove that any split-table prediction induces a
feasible parenthesization, so structural models can never report a cost below the
DP optimum. On a held-out benchmark of 500 synthetic chains ($n{=}5$--$50$) drawn
from four dimension distributions, GraphMCMNet attains 0.029\% MAPE, a
log-scale $R^2$ of 0.999999, and reproduces the exact DP split table in 93\% of
cases. PointerMCMNet also demonstrates exceptional accuracy with 0.148\% MAPE. Crucially, both structural models maintain 100\% mathematical validity. In contrast, the tree
ensembles---despite extensive feature engineering---predict costs below the DP
lower bound in 78--82\% of cases. Wilcoxon signed-rank tests confirm all pairwise
differences ($p{<}0.001$), with very large effect sizes ($|d|{>}1.0$) between the
structural and statistical families. Trained only on chains of length
$n{\le}35$, GraphMCMNet achieves its \emph{lowest} error (0.019\% MAPE) on
out-of-distribution chains of length $41$--$50$. Preliminary evaluation on
matrix chains extracted from real neural architectures (ResNet-50, BERT,
GPT-2, ViT) corroborates these findings. On our benchmark,
these results indicate that predicting combinatorial structure, rather than
regressing unconstrained costs, is the decisive design choice for learned MCM
approximation.
\end{abstract}

\begin{IEEEkeywords}
matrix chain multiplication, graph neural network, neural combinatorial optimization,
pointer network, dynamic programming, neural algorithmic reasoning, structural reasoning
\end{IEEEkeywords}

\section{Introduction}
The Matrix Chain Multiplication (MCM) problem is a cornerstone of dynamic
programming: given $n$ matrices, find the parenthesization that minimizes the
total number of scalar multiplications. The classical DP solution runs in $O(n^3)$
time and $O(n^2)$ space~\cite{b4}. Hu and Shing~\cite{b2} later proposed an
$O(n\log n)$ algorithm; despite its asymptotic advantage, it is rarely deployed
in practice due to implementation complexity~\cite{b3}. Exact optimization thus
remains a bottleneck in deep learning compilers~\cite{b23,b24}, database query
optimizers~\cite{b17,b19}, and tensor contraction engines that must evaluate
thousands of chains repeatedly.

Prior learned approaches have relied on heuristic or evolutionary
search~\cite{b5}, which offer neither optimality nor feasibility guarantees.
Recent advances in neural combinatorial optimization~\cite{b7,b8}, GNNs for
structured problems~\cite{b20,b21,b22}, and neural algorithmic
reasoning~\cite{b25,b26} suggest that architectures respecting a problem's
combinatorial structure can approximate exact solvers with high fidelity.

This paper makes the following contributions:
\begin{enumerate}
    \item We introduce \textbf{GraphMCMNet}, a graph neural network that maps
    the DP sub-problem lattice to an explicit graph and predicts optimal split
    points, attaining 0.029\% MAPE. We also design \textbf{PointerMCMNet}, an attention-based sequence model that achieves 0.148\% MAPE. Both structural models attain 100\% mathematical validity on our 500-chain held-out benchmark.
    \item We formally prove (Theorem~\ref{thm:validity}) that any model
    predicting per-sub-chain split points induces a feasible parenthesization,
    so its reported cost can never fall below the DP optimum---a guarantee that
    no unconstrained cost regressor can offer.
    \item We provide a controlled comparison of structural reasoning (GNN and
    Pointer Network) versus statistical estimation (XGBoost and Random Forest
    with 145 engineered features), observing a three-order-of-magnitude accuracy
    gap on our benchmark, with all differences statistically significant
    ($p{<}0.001$, $|d|{>}1.0$).
    \item We demonstrate length extrapolation: trained only on $n{\le}35$,
    GraphMCMNet achieves its lowest error (0.019\% MAPE) on out-of-distribution
    chains 43\% longer than any training instance.
\end{enumerate}

We deliberately restrict our claims to the synthetic benchmark studied here;
Section~\ref{sec:limitations} discusses external validity and the experiments
required to extend these conclusions to production workloads.

\section{Related Work}
\textbf{Classical MCM Algorithms.} Godbole~\cite{b1} formalized the $O(n^3)$
dynamic program. Hu and Shing~\cite{b2} achieved $O(n\log n)$ via a reduction
to convex polygon partitioning. L\'{o}pez et al.~\cite{b3} recently showed that
few parenthesizations are ever optimal. Despite its asymptotic advantage,
Hu--Shing is rarely deployed in practical optimizers.

\textbf{Neural Combinatorial Optimization.} Vinyals et al.~\cite{b7} introduced
Pointer Networks for sequence-to-sequence combinatorial problems. Bello
et al.~\cite{b8} extended them with reinforcement learning. Bahdanau
et al.~\cite{b9} proposed additive attention underpinning pointer decoders.
These approaches lack explicit structural inductive biases.

\textbf{GNNs for Optimization.} Cappart et al.~\cite{b20} surveyed GNNs in
combinatorial optimization. Bengio et al.~\cite{b21} provided a methodological
framework for ML-driven optimization. Vesselinova et al.~\cite{b22} demonstrated
that learning on graphs improves generalization for combinatorial tasks. Our work
extends this line by constructing the DP sub-problem graph as an explicit GNN input.

\textbf{Neural Algorithmic Reasoning.} Veli\v{c}kovi\'{c} and
Blundell~\cite{b25} articulated the neural algorithmic reasoning blueprint; the
CLRS benchmark~\cite{b26} showed that message-passing networks aligned with an
algorithm's intermediate state generalize far better than unaligned models.
Ibarz et al.~\cite{b27} scaled this to 30 algorithms in the CLRS-30
benchmark, and Mahdavi et al.~\cite{b28} studied out-of-distribution
generalization in neural algorithmic reasoning---directly relevant to our
length-extrapolation experiment. Bevilacqua et al.~\cite{b29} recently
showed that causal regularisation further improves algorithmic learners.
Our split-table prediction task is an instance of this paradigm applied to the MCM
dynamic program.

\textbf{Learned Cost Models.} Neo~\cite{b17} and Balsa~\cite{b19} train neural
models for query optimization. TVM~\cite{b23} and Tiramisu~\cite{b24} use learned
cost models in compiler optimization. These approaches regress unconstrained
scalars, offering no feasibility guarantee---a limitation our structural approach
addresses directly.

\section{Problem Formulation}
\label{sec:formulation}

Given a chain of $n$ matrices $A_1, A_2, \ldots, A_n$, where matrix $A_i$ has dimensions $d_{i-1} \times d_i$, the Matrix Chain Multiplication (MCM) problem seeks a parenthesization that minimizes the total number of scalar multiplications. Although the final product is identical regardless of evaluation order, the cost varies dramatically: for a 5-matrix chain, one parenthesization may require 10,500 multiplications while another demands over 200,000. The classical dynamic programming (DP) recurrence~\cite{b4} solves this exactly:
\begin{equation}
\label{eq:dp}
m[i][j] = \begin{cases}
0 & \text{if } i = j \\
\displaystyle\min_{i \leq k < j} \bigl\{ m[i][k] + m[k{+}1][j] + d_{i-1} \cdot d_k \cdot d_j \bigr\} & \text{if } i < j
\end{cases}
\end{equation}

Here $m[i][j]$ is the minimum cost for sub-chain $A_i \cdots A_j$, and the optimal split table $s[i][j] = \arg\min_k\{\cdot\}$ records the best split point for each sub-chain. Computing the full solution requires $O(n^3)$ time and $O(n^2)$ space---a severe bottleneck for real-time systems such as deep learning compilers and query optimizers that must evaluate thousands of chains per second.

\medskip
\noindent\textbf{Central Hypothesis.}\quad Can machine learning models approximate
this cubic-time solver while preserving solution quality and feasibility?
We investigate two fundamentally different framings:

\begin{enumerate}
\item \textbf{Structural Reasoning (Split Prediction):} The model predicts the
split table $\hat{s}[i][j]$ for every sub-chain. The cost is then \emph{computed
exactly} from these splits. Because every predicted split corresponds to a real
parenthesization, the resulting cost is always \emph{mathematically valid}:
$\hat{c} \geq m[1][n]$ by construction.

\item \textbf{Statistical Estimation (Cost Regression):} The model directly
regresses a scalar cost $\hat{c}$ from hand-crafted features. No constraint
enforces that $\hat{c}$ corresponds to any valid parenthesization---so the model
may predict $\hat{c} < m[1][n]$, a cost that is \emph{mathematically infeasible}.
\end{enumerate}

The following theorem formalizes why the first framing is safe by construction.

\begin{theorem}[Structural Validity Guarantee]
\label{thm:validity}
Let $\hat{s}$ be any predicted split table such that
$\hat{s}[i][j] \in \{i, i{+}1, \ldots, j{-}1\}$ for every sub-chain
$1 \le i < j \le n$. Define the induced cost recursively:
\begin{equation}
\hat{m}[i][j] =
\begin{cases}
0 & i = j\\
\hat{m}[i][k^*] + \hat{m}[k^*{+}1][j] + d_{i-1} d_{k^*} d_j & i < j
\end{cases}
\end{equation}
with $k^* = \hat{s}[i][j]$. Then the evaluation order induced by $\hat{s}$ is a
legal full parenthesization of $A_i \cdots A_j$ for every $(i,j)$, and
$\hat{m}[1][n] \ge m[1][n]$, with equality if and only if $\hat{s}$ agrees with
some optimal split table on every sub-chain reachable from $(1,n)$.
\end{theorem}

\begin{proof}
By strong induction on sub-chain length $\ell = j - i + 1$. Base case
($\ell{=}1$): trivial. Inductive step: for $(i,j)$ let $k^* = \hat{s}[i][j]$.
By hypothesis $i \le k^* < j$, so $(i,k^*)$ and $(k^*{+}1,j)$ are well-formed
and, by induction, each admits a legal parenthesization. Composing them gives a
legal parenthesization of $A_i \cdots A_j$ with cost $\hat{m}[i][j]$. Since
$m[i][j]$ minimizes over all legal parenthesizations, $\hat{m}[i][j] \ge m[i][j]$.
Applying this at $(1,n)$ yields $\hat{m}[1][n] \ge m[1][n]$.
\end{proof}

\begin{remark}
Theorem~\ref{thm:validity} holds for \emph{any} predictor whose output is a split
table, including an untrained one: the guarantee is architectural, not learned. An
unconstrained regressor may output $\hat{c} < m[1][n]$, a cost achieved by no
parenthesization. We measure the fraction of predictions satisfying
$\hat{c} \ge m[1][n]$ as \emph{mathematical validity} in Section~V.
\end{remark}

This distinction---predicting \emph{combinatorial structure} versus
\emph{unconstrained numbers}---is the central axis of this study.


%% ================================================================
%%   SECTION IV: PROPOSED METHODOLOGY
%% ================================================================

\section{Proposed Methodology}
\label{sec:methodology}

\subsection{System Overview}
\label{sec:overview}


Fig.~\ref{fig:architecture} presents the overall framework. Given a dimension sequence $[d_0, d_1, \ldots, d_n]$, two paradigms operate in parallel. The \emph{Structural Reasoning} branch deploys GraphMCMNet and PointerMCMNet to predict split positions and compute exact costs. The \emph{Statistical Estimation} branch uses XGBoost and Random Forest to regress costs from engineered features. All predictions are evaluated against identical DP ground truth.

\begin{figure}[htbp]
\centering
\includegraphics[width=\columnwidth]{fig_system_architecture.png}
\caption{System architecture. Structural models predict split tables and compute costs exactly; statistical models regress scalar costs from 145 engineered features. Both are evaluated against DP ground truth.}
\label{fig:architecture}
\end{figure}

\subsection{GraphMCMNet: Learning the DP Structure}
\label{sec:gnn}

\begin{figure*}[t]
\centering
\includegraphics[width=\textwidth]{neuralnetwork.jpeg}
\caption{Neural architectures for structural reasoning. (a) PointerMCMNet combines a 6-layer Transformer encoder with a Bahdanau attention pointer decoder and structural validity mask. (b) GraphMCMNet maps the MCM dynamic programming problem to a sub-problem graph, refines embeddings via 6 gated message-passing layers, and scores splits using parent-child triples.}
\label{fig:neural_architectures}
\end{figure*}

The key insight behind GraphMCMNet is that the DP recurrence in~(\ref{eq:dp}) naturally defines a graph over sub-problems---and a graph neural network can learn to traverse it.

\medskip
\noindent\textbf{Sub-Problem Graph Construction.}\quad Each sub-chain $(i,j)$ for $1 \leq i \leq j \leq n$ becomes a node, yielding $n(n{+}1)/2$ nodes total (e.g., 10 nodes for $n{=}4$, arranged in a triangular structure from leaves to root). Three types of edges connect these nodes: (i)~\emph{DP child edges} that mirror the recurrence---connecting each parent $(i,j)$ bidirectionally to its left child $(i,k)$ and right child $(k{+}1,j)$ for every valid split $k$; (ii)~\emph{same-length neighbor edges} linking adjacent sub-chains of identical length, enabling lateral information flow; and (iii)~implicit \emph{self-loops} through gated residual connections. This construction gives each node direct access to every sub-problem it depends on in the DP table.

Each node receives a compact 10-dimensional feature vector capturing its local context: log-transformed boundary dimensions, their product, normalized sub-chain length, mean and standard deviation of log-dimensions, min/max statistics, the position of the minimum dimension, and a leaf indicator. These features are designed to be informative yet lightweight, as the heavy lifting is delegated to message passing.

\medskip
\noindent\textbf{Architecture.}\quad GraphMCMNet (Fig.~\ref{fig:neural_architectures}(b)) processes this graph in four stages. First, raw features are projected into a 128-dimensional latent space via a linear layer with LayerNorm and SiLU activation. Next, $K{=}6$ rounds of \emph{gated message passing} iteratively refine the embeddings: at each layer, every node collects messages from its neighbors via a two-layer MLP, mean-aggregates them, and updates its own embedding through a \emph{learned gate} that controls how much new information to incorporate. This gating mechanism---a sigmoid-activated linear layer that element-wise modulates the update---prevents over-smoothing and lets the network selectively integrate structural information while preserving local features. Each layer concludes with LayerNorm and a residual connection.

\medskip
\noindent\textbf{Parameter Selection.}\quad The embedding dimension $d{=}128$ was selected via grid search on the validation set: $d{=}64$ reduced exact-match rate by 4.2\%, while $d{=}256$ yielded negligible improvement at double the memory cost. $K{=}6$ message-passing rounds approximately match half the diameter of the DP sub-problem graph for $n{\le}50$ (diameter ${\le}2\log_2 n {\approx} 12$), ensuring sufficient information propagation. The 6-layer Transformer encoder in PointerMCMNet mirrors this depth for controlled comparison.

After message passing, the model must decide: for each sub-chain $(i,j)$, which split point $k$ is optimal? This is answered by the \textbf{split scorer}, which concatenates the embeddings of the parent node $(i,j)$ with those of the candidate children $(i,k)$ and $(k{+}1,j)$, and feeds the triple through a three-layer MLP:
\begin{equation}
\text{score}(k \mid i,j) = \text{MLP}_{\text{split}}\bigl(\mathbf{h}_{(i,j)} \;\|\; \mathbf{h}_{(i,k)} \;\|\; \mathbf{h}_{(k+1,j)}\bigr)
\label{eq:split_score}
\end{equation}

The predicted split is simply $\hat{s}[i][j] = \arg\max_k \, \text{score}(k \mid i,j)$. Crucially, this scoring is fully vectorized: for each sub-chain length $L$, all sub-chains and their candidates are scored in a single batched forward pass. An auxiliary cost head predicts $\log(1{+}m[1][n])$ from the root node embedding, providing a direct gradient signal for cost accuracy during training.


\subsection{PointerMCMNet: Sequential Split Selection}
\label{sec:pointer}

While GraphMCMNet explicitly models the DP graph, PointerMCMNet (Fig.~\ref{fig:neural_architectures}(a)) takes a different approach: it processes the dimension sequence left-to-right using a Transformer encoder~\cite{b11} and ``points'' to split positions using attention~\cite{b7,b9}. This tests whether a sequential encoder, without explicit graph structure, can reconstruct the DP relationships implicitly.

\medskip
\noindent\textbf{Encoder.}\quad Each of the $n{+}1$ dimension positions receives an 8-dimensional feature vector: the log-dimension, products with left and right neighbors, the triple product, normalized position, normalized magnitude, and binary indicators for local valleys and peaks. These features are projected to 128 dimensions, enriched with sinusoidal positional encodings, and processed by a 6-layer Transformer encoder (8 heads, feed-forward dimension 512, GELU activation, dropout 0.1).

\medskip
\noindent\textbf{Pointer Decoder.}\quad For each sub-chain $(i,j)$ of length $L \geq 3$, a query is formed by concatenating the boundary hidden states $\mathbf{h}_{i-1}$ and $\mathbf{h}_j$ with a 32-dimensional learnable length embedding, then projecting through an MLP. Each candidate split position $k$ is scored using Bahdanau (additive) attention: $\text{score}(k \mid i,j) = \mathbf{v}^T \tanh(\mathbf{W}_1 \mathbf{q}_{(i,j)} + \mathbf{W}_2 \mathbf{h}_k)$. The predicted split is $\hat{s}[i][j] = i + \arg\max_k \, \text{softmax}(\text{scores})$, and an auxiliary cost head (masked mean pooling followed by an MLP) provides the same training signal as in GraphMCMNet.


\subsection{Statistical Baselines}
\label{sec:baselines}

To rigorously test whether structural reasoning is truly necessary, we implement strong tree-ensemble baselines---XGBoost~\cite{b12} and Random Forest~\cite{b13}---armed with an extensive 145-dimensional feature vector. If exhaustive feature engineering could solve MCM, these models would demonstrate it.

The 145 features span four groups: \emph{global statistics} (30 features: chain length, dimension min/max/mean/std/median, percentiles, pattern indicators), \emph{positional features} (8 features: locations of extrema, quartile means, trend slope), \emph{cost proxies} (6 features: log-costs from four greedy heuristics---left-to-right, right-to-left, minimum-first, and balanced splitting), and \emph{padded sequence features} (101 features: the full dimension sequence and pairwise products, log-transformed and zero-padded).

Both XGBoost and Random Forest use a \emph{dual-head ensemble}: a \emph{direct head} that predicts $\log(1{+}c)$ and a \emph{ratio head} that predicts the ratio $c / c_{\text{greedy}}$. Final predictions blend both: $\hat{c} = 0.3 \times \hat{c}_{\text{direct}} + 0.7 \times c_{\text{greedy}} \times \hat{r}$, anchoring predictions to a known-good greedy reference.


\subsection{Training Strategy}
\label{sec:training}

\noindent\textbf{Dataset.}\quad 120,000 matrix chains are generated with $n \in [5, 50]$, balanced across four dimension distributions: \emph{Uniform} ($d_i \sim U[10,500]$), \emph{Spiky} (alternating extremes), \emph{Bottleneck} (near-zero inner dimensions), and \emph{Monotone} (linearly increasing). Each sample stores the dimension sequence, DP-optimal cost, and the complete split table. The data is split 72\%/13\%/15\% for training, validation, and testing.

\medskip
\noindent\textbf{Curriculum Learning.}\quad Both neural models are trained with a multi-stage curriculum (Fig.~\ref{fig:training}, Table~\ref{tab:curriculum}) that progressively increases chain length, mirroring how a student masters short problems before tackling longer ones.

\begin{figure}[htbp]
\centering
\includegraphics[width=\columnwidth]{fig_training_pipeline.png}
\caption{Four-stage curriculum learning pipeline. Chain length increases progressively; the GNN stops at Stage~3 ($n{\leq}35$) to test out-of-distribution generalization.}
\label{fig:training}
\end{figure}

\begin{table}[htbp]
\caption{Curriculum Learning Schedule}
\label{tab:curriculum}
\begin{center}
\renewcommand{\arraystretch}{1.2}
\setlength{\tabcolsep}{6pt}
\begin{tabular}{|c|c|c|c|c|c|}
\hline
\textbf{Stage} & \textbf{$n_{\max}$} & \textbf{Epochs} & \textbf{LR} & \textbf{Batch} & \textbf{Mode} \\
\hline
1 & 10  & 30 & $10^{-3}$   & 128 & Cached \\
2 & 20  & 30 & $5{\times}10^{-4}$ & 96  & Cached \\
3 & 35  & 5  & $2{\times}10^{-4}$ & 48  & Low-mem \\
4 & 50  & 8  & $10^{-4}$   & 32  & Low-mem \\
\hline
\end{tabular}
\end{center}
\end{table}

A crucial design decision: the GNN is trained \emph{only through Stage~3} ($n_{\max}{=}35$), deliberately withholding all $n{>}35$ chains. This sets up a rigorous test of out-of-distribution generalization---can the model compose solutions for chains 43\% longer than anything it has seen?

\medskip
\noindent\textbf{Loss Function.}\quad Training optimizes a composite objective that balances structural accuracy with cost fidelity:
\begin{equation}
\mathcal{L} = \underbrace{0.9 \cdot \sum_{L=2}^{n} \frac{L \cdot \text{CE}_L}{\sum_{L'} L' \cdot |\mathcal{S}_{L'}|}}_{\text{length-weighted split loss}} + \underbrace{0.1 \cdot \text{LogCosh}\bigl(\hat{y}_{\text{cost}} - y_{\text{cost}}\bigr)}_{\text{auxiliary cost loss}}
\label{eq:loss}
\end{equation}
The split loss uses cross-entropy weighted by sub-chain length $L$, so longer (harder) sub-chains receive proportionally more gradient signal. The auxiliary cost loss uses LogCosh---a smooth, outlier-robust alternative to MSE---to encourage accurate cost estimates without dominating the structural objective.

\medskip
\noindent\textbf{Optimization.}\quad AdamW optimizer ($\beta_1{=}0.9$, $\beta_2{=}0.98$, weight decay $10^{-4}$) with cosine annealing warm restarts, gradient clipping at norm 1.0, and mixed-precision training (AMP) on an NVIDIA RTX~4050. Early stopping (patience~15) prevents overfitting. Total training: ${\sim}$25 GPU-hours.



\subsection{Evaluation Protocol}
\label{sec:evalprotocol}

A held-out set of 500 matrix chains (seed${=}$42, $n{=}5$--$50$) is balanced across all four distributions (125 each). Five metrics capture complementary quality dimensions: \emph{MAPE} (average cost accuracy), \emph{$R^2$ on log-scale} (correlation across magnitudes), \emph{mathematical validity} (fraction of predictions where $\hat{c} \geq m[1][n]$---any violation means the model hallucinated a non-existent parenthesization), \emph{exact match} (predicted splits $\equiv$ DP splits, neural models only), and \emph{F1-score} at a 1\% error threshold. All pairwise model comparisons use Wilcoxon signed-rank tests with Cohen's $d$ for effect size.

\subsection{Reproducibility}
\label{sec:reproducibility}
All experiments use a fixed random seed (42) for data generation and evaluation. Training seeds are set via PyTorch, NumPy, and Python's \texttt{random} module to ensure deterministic initialization. The 500-chain evaluation set is fully specified by its seed and distribution parameters. Source code, trained checkpoints, and evaluation scripts are publicly available at \url{https://github.com/abhilash0042/ml-matrix-chain-optimizer}. All reported metrics can be reproduced with a single evaluation script invocation.

%% ================================================================
%%   SECTION V: RESULTS AND DISCUSSION
%% ================================================================

\section{Results and Discussion}

\subsection{Experimental Setup}
\label{sec:setup}

Evaluation uses 500 held-out matrix chains (seed$=$42, $n{=}5$--$50$) balanced across four dimension distributions: \emph{Uniform} ($d_i{\sim}U[10,500]$), \emph{Spiky} (alternating extreme scales), \emph{Bottleneck} (one near-zero inner dimension), and \emph{Monotone} (linearly increasing). Ground-truth optimal costs and split tables come from exact DP. The four models under comparison are \textbf{GraphMCMNet} (6-layer graph attention, 3.7\,M parameters, trained on $n{\leq}35$ via 3-stage curriculum), \textbf{PointerMCMNet} (6-layer transformer encoder with Bahdanau decoder, 5.4\,M parameters, trained to $n{\leq}50$), and two tree-ensemble baselines---\textbf{XGBoost} and \textbf{Random Forest}---fed 145 hand-crafted features (greedy cost proxies, spectral properties, dimension ratios). Neural training ran on an NVIDIA RTX~4050 for roughly 25 GPU-hours. Metrics include MAPE, log-scale $R^2$, mathematical validity (predicted cost ${\geq}$ DP optimum), exact split-table match rate, and F1-score at 1\% error.

\subsection{Overall Performance}
\label{sec:mcmresults}

To ground the numbers in the actual problem: given a 5-matrix chain with dimensions $[10, 100, 5, 50, 10, 30]$, DP finds the optimum at \textbf{10,500} multiplications via $((A_1 A_2)((A_3 A_4) A_5))$ (Fig.~\ref{fig:worked_example}). GraphMCMNet recovers this exact split table. XGBoost predicts 18,347---a 74.7\% overshoot. On a 40-matrix spiky chain the GNN lands within 0.09\% of DP; Random Forest overshoots by over 200\%. Across all 500 chains the GNN reproduces the \emph{exact} DP split table 93\% of the time (Table~\ref{tab:overall}).

\begin{figure}[htbp]
\centering
\includegraphics[width=\columnwidth]{workedexample.jpeg}
\caption{DP sub-problem tree for a 5-matrix chain with $d{=}[10,100,5,50,10,30]$. Each internal node shows the sub-chain $(i,j)$, optimal split $k$, and minimum cost. The root $(1,5)$ with $k{=}2$ yields the parenthesization $((A_1 A_2)((A_3 A_4)A_5))$ at cost 10,500. GraphMCMNet recovers this exact tree.}
\label{fig:worked_example}
\end{figure}

\medskip
\begin{table}[htbp]
\caption{Overall Model Performance on 500 Test Chains ($n{=}5$--$50$). Best in \textbf{bold}. MAPE includes 95\% bootstrap CI.}
\label{tab:overall}
\begin{center}
\renewcommand{\arraystretch}{1.3}
\setlength{\tabcolsep}{5pt}
\begin{tabular}{|l|c|c|c|c|c|}
\hline
\textbf{Model} & \textbf{MAPE} & \textbf{$R^2$} & \textbf{Valid} & \textbf{Exact} & \textbf{F1} \\
               & \textbf{(\%)} & \textbf{(log)} & \textbf{(\%)}  & \textbf{Match} & \textbf{(\%)} \\
\hline
GraphMCMNet   & \textbf{0.029} \scriptsize{[0.01, 0.05]} & \textbf{0.999999} & \textbf{100.0} & \textbf{93.0\%} & \textbf{99.60} \\
PointerMCMNet & 0.148 \scriptsize{[0.08, 0.23]}          & 0.9999          & \textbf{100.0} & 88.8\%          & 98.48          \\
XGBoost       & 36.187         & 0.9067          & 22.0           & ---             & ---            \\
Random Forest & 43.284         & 0.8922          & 18.2           & ---             & ---            \\
\hline
\end{tabular}
\end{center}
\end{table}

\medskip
The GNN deviates from the DP cost by fewer than 3 parts in 10,000. Both neural models maintain 100\% \emph{mathematical validity}: because each predicts discrete split positions rather than a raw scalar, the resulting cost always corresponds to some real parenthesization and can never dip below the DP lower bound. Tree ensembles, by contrast, violate that bound in 78--82\% of cases. The root cause is architectural: gradient-boosted trees and random forests treat the cost as an ordinary continuous target with no notion of the recursive feasibility that DP enforces. The GNN holds a $5.1{\times}$ MAPE edge over the Pointer Network (0.029\% vs.\ 0.148\%) because it encodes the $O(n^2)$ sub-problem dependencies as explicit graph edges; the Pointer Network must reconstruct those relationships implicitly through self-attention. However, PointerMCMNet remains an exceptionally strong baseline: it attains 100\% mathematical validity and a near-zero MAPE of 0.148\% without any explicit structural biases, demonstrating the power of sequence-to-sequence models for combinatorial tasks. While the GNN's graph structure provides the final edge, both structural paradigms comprehensively dominate the statistical ensembles.

These validity numbers carry practical weight (Fig.~\ref{fig:validity}). A cost model inside a compiler or query optimizer that routinely returns impossible values would steer scheduling decisions toward multiplication orders that cannot physically exist, corrupting every downstream pass that relies on the estimate. The 100\% validity of the structural models eliminates this failure mode entirely.

\begin{figure}[htbp]
\centering
\includegraphics[width=\columnwidth]{fig2_validity_accuracy.png}
\caption{Model Validity vs. Accuracy. Only structural models (GraphMCMNet and PointerMCMNet) maintain 100\% mathematical validity.}
\label{fig:validity}
\end{figure}


\subsection{Generalization and Robustness}
\label{sec:scale}

Table~\ref{tab:breakdown} breaks MAPE down by chain length and dimension distribution. The error distribution is further visualized in Fig.~\ref{fig:distribution}.

\medskip
\begin{table}[htbp]
\caption{MAPE (\%) by Chain Length and Distribution. $^\dagger$OOD: GNN trained only to $n{=}35$.}
\label{tab:breakdown}
\begin{center}
\renewcommand{\arraystretch}{1.3}
\setlength{\tabcolsep}{8pt}
\begin{tabular}{|l|c|c|c|c|}
\hline
\textbf{Condition} & \textbf{GNN} & \textbf{Pointer} & \textbf{XGBoost} & \textbf{RF} \\
\hline
\multicolumn{5}{|l|}{\textit{By chain length}} \\
\hline
$n{=}5$--$20$             & \textbf{0.027} & 0.092 & 33.09 & 45.22 \\
$n{=}21$--$40$            & \textbf{0.036} & 0.082 & 39.57 & 43.37 \\
$n{=}41$--$50^\dagger$    & \textbf{0.019} & 0.391 & 35.62 & 39.23 \\
\hline
\multicolumn{5}{|l|}{\textit{By dimension distribution}} \\
\hline
Uniform    & \textbf{0.003} & 0.032 & 54.18 & 51.18 \\
Spiky      & \textbf{0.114} & 0.500 & 41.40 & 46.68 \\
Bottleneck & \textbf{0.000} & 0.059 & 41.97 & 62.94 \\
Monotone   & \textbf{0.000} & \textbf{0.000} & 7.19 & 12.34 \\
\hline
\end{tabular}
\end{center}
\vspace{1mm}
{\footnotesize $^\dagger$GNN trained on $n{\leq}35$ only, evaluated beyond its training horizon.}
\end{table}

\begin{figure}[htbp]
\centering
\includegraphics[width=\columnwidth]{fig3_mape_distribution.png}
\caption{Error distribution across different matrix chain properties. The GNN consistently achieves near-zero MAPE.}
\label{fig:distribution}
\end{figure}
\medskip
The GNN was never shown chains longer than 35 during training, yet it records its
\emph{lowest} error (0.019\%) on the $n{=}41$--$50$ bucket---better than its own
in-distribution average (0.032\%). Longer chains are composed of shorter sub-chains
the model has already mastered; the graph edges let it assemble solutions for larger
instances out of familiar pieces, consistent with findings in neural algorithmic
reasoning that algorithm-aligned architectures extrapolate in size~\cite{b26}. This sustained structural accuracy on out-of-distribution lengths is highlighted in Fig.~\ref{fig:scaling}.

\begin{figure}[htbp]
\centering
\includegraphics[width=\columnwidth]{fig4_structural_scaling.png}
\caption{Exact Match Rate vs. chain length. GraphMCMNet sustains high structural accuracy even for out-of-distribution lengths ($n>35$).}
\label{fig:scaling}
\end{figure}
The Pointer Network, lacking explicit structure, degrades $4.3{\times}$ on the same
length bucket (0.391\% vs.\ 0.092\%), though we note it was trained on the full
length range, so this comparison measures architectural robustness rather than strict
OOD extrapolation. Tree models hover between 31\% and 47\% error at every length,
reinforcing that the gap is not about scale but about the mismatch between
unconstrained regression and combinatorial structure.

Bottleneck and Monotone chains are solved essentially perfectly by the GNN
(0.000\% MAPE). Spiky chains are the hardest for every model, but the GNN's worst
case there (0.114\%) is still smaller than XGBoost's best case on any distribution
(7.19\% on Monotone).

\subsection{Statistical Significance}
\label{sec:stats}


Table~\ref{tab:wilcoxon} reports the full significance analysis. All six pairwise
comparisons are significant at $p{<}0.001$ and survive Bonferroni correction
($\alpha=0.05/6\approx0.0083$). Structural-vs-statistical comparisons exhibit very
large effect sizes ($|d|{>}1.0$); GNN-vs-Pointer shows a small but significant
effect ($d{=}0.13$). In head-to-head win rates, the GNN wins 100\% of chains
against both tree baselines and 91\% against the Pointer Network.

\begin{table}[htbp]
\caption{Wilcoxon Signed-Rank Tests on Per-Chain Errors (Two-Sided,
Bonferroni-corrected threshold $\alpha=0.0083$).}
\label{tab:wilcoxon}
\centering
\setlength{\tabcolsep}{4pt}
\begin{tabular}{@{}lccc@{}}
\toprule
Comparison & $W$ & $p$-value & Cohen's $d$ \\
\midrule
Pointer vs.\ GNN     & 302      & $1.5{\times}10^{-4}$  & $0.13$ \\
Pointer vs.\ XGBoost & 3        & $1.3{\times}10^{-83}$ & $-1.34$ \\
Pointer vs.\ RF      & 11       & $1.4{\times}10^{-83}$ & $-1.09$ \\
GNN vs.\ XGBoost     & 60       & $1.8{\times}10^{-83}$ & $-1.34$ \\
GNN vs.\ RF          & 0        & $1.3{\times}10^{-83}$ & $-1.09$ \\
XGBoost vs.\ RF      & 44{,}082 & $9.7{\times}10^{-9}$  & $-0.25$ \\
\bottomrule
\end{tabular}
\end{table}


\subsection{Evaluation on Real Neural Architectures}
\label{sec:realarch}

To evaluate whether our findings transfer from synthetic distributions to production workloads, we extracted 98 distinct matrix chains representing the exact tensor dimensions from widely deployed neural architectures (including ResNet-18/34/50/101/152, BERT-base/large, GPT-2, ViT, EfficientNet, MobileNet, DenseNet, and T5). For instance, the sequence of projections in a BERT-base transformer block induces a weight matrix chain of dimensions $d{=}[768, 768, 768, 3072, 768]$.

As shown in Table~\ref{tab:realarch}, both structural neural models (GraphMCMNet and PointerMCMNet) generalize perfectly to these real-world chains, achieving 0.000\% MAPE, 100\% mathematical validity, and recovering the exact DP optimum on all 98 instances. This seemingly flawless performance is expected: real architecture chains tend to be short ($n \le 8$) and lack the adversarial, spiky dimensional patterns that challenged both GraphMCMNet and PointerMCMNet in our synthetic evaluation. By contrast, XGBoost averages over 51.0\% error and Random Forest exceeds 58\% error, with neither producing a mathematically valid prediction. This underscores that statistical models fail to capture the underlying combinatorial constraints even on simpler, well-behaved distributions.

\begin{table}[htbp]
\caption{Benchmark on Matrix Chains Extracted from Real Neural Architectures. Metrics aggregated over 98 chains.}
\label{tab:realarch}
\begin{center}
\renewcommand{\arraystretch}{1.3}
\setlength{\tabcolsep}{8pt}
\begin{tabular}{@{}lccc@{}}
\toprule
\textbf{Model} & \textbf{MAPE} & \textbf{Valid (\%)} & \textbf{Exact Match (\%)} \\
\midrule
GraphMCMNet   & \textbf{0.000\%} & \textbf{100.0} & \textbf{100.0} \\
PointerMCMNet & \textbf{0.000\%} & \textbf{100.0} & \textbf{100.0} \\
XGBoost       & 51.017\%         & 0.0            & 0.0            \\
Random Forest & 58.420\%         & 0.0            & 0.0            \\
\bottomrule
\end{tabular}
\end{center}
\end{table}

\subsection{Latency and Ablation}
\label{sec:sigablation}

Table~\ref{tab:complexity} lists computational complexity and measured speedups. Ablation conclusions draw directly on Tables~\ref{tab:overall} and~\ref{tab:breakdown}.

\medskip
\begin{table}[htbp]
\caption{Computational Complexity and Inference Latency}
\label{tab:complexity}
\begin{center}
\renewcommand{\arraystretch}{1.3}
\setlength{\tabcolsep}{5pt}
\begin{tabular}{@{}lcccc@{}}
\toprule
\textbf{Model} & \textbf{Complexity} & \textbf{Params} & \multicolumn{2}{c}{\textbf{Latency (ms/chain)$^\ast$}} \\
 & & & $n{=}10$ & $n{=}50$ \\
\midrule
DP (Exact)       & $O(n^3)$      & --- & 0.05 & 3.97 \\
Greedy (Best-4)  & $O(n^2)$      & --- & 0.02 & 0.18 \\
Hu--Shing$^\S$   & $O(n\log n)$  & --- & \multicolumn{2}{c}{same output as DP} \\
GraphMCMNet      & $O(n^2 K)$    & 3.7\,M & 27.3 & 349$^\dagger$ \\
PointerMCMNet    & $O(n^2)$      & 5.4\,M & 28.8 & 225 \\
XGBoost / RF     & $O(n^2)^\ddagger$ & --- & 12.3 & 17.1 \\
\bottomrule
\end{tabular}
\end{center}
\vspace{1mm}
{\footnotesize $^\ast$Single-chain CPU inference, mean over 100 runs
(warm-up: 10). $^\dagger$High variance due to graph memory allocation.
$^\ddagger$Feature extraction dominates. $K{=}6$ GNN layers.
$^\S$Hu--Shing computes the identical optimal cost as DP; see text.
GPU batching yields ${>}100{\times}$ speedup for all learned models.}
\end{table}

\medskip
We emphasize that neural inference is \emph{not} constant time: it scales as
$O(n^2 K)$ for the GNN (with fixed message-passing rounds $K$) and $O(n^2)$ for
CPU latency. On single chains, Python DP (0.05--3.97\,ms) is actually faster than
unoptimized neural inference (27--349\,ms) because neural models pay a fixed
per-chain graph-construction and device-transfer overhead. The practical advantage
emerges under \emph{GPU batching}: when thousands of chains are evaluated
concurrently on accelerator hardware, learned models achieve ${>}100{\times}$
aggregate throughput---the target regime for JIT compilers~\cite{b23} and query
optimizers~\cite{b17}. We note that the Hu--Shing algorithm~\cite{b2} computes the \emph{same} exact optimal cost as $O(n^3)$ DP, differing only in runtime complexity ($O(n\log n)$). Since our research question concerns approximation \emph{quality}---whether learned models can match the exact optimum while preserving structural validity---the accuracy and validity conclusions in Tables~\ref{tab:overall}--\ref{tab:breakdown} hold identically regardless of which exact solver serves as ground truth. For single-chain serial latency at large~$n$, Hu--Shing would outperform all methods; our target regime is batch inference on GPU, where learned models dominate.

Three ablation observations stand out. \emph{(i)}~Swapping the graph encoder for
the sequential one raises MAPE $5.1{\times}$ (0.029\%$\to$0.148\%), confirming
that explicit DP edges matter. \emph{(ii)}~Replacing split prediction with direct
cost regression drops validity from 100\% to 22\%; the structural guarantee is
architectural and cannot be post-hoc recovered. \emph{(iii)}~OOD MAPE
\emph{improves} versus in-distribution (0.019\% vs.\ 0.032\%), evidence that the
GNN has internalized the compositional structure of the DP recurrence rather than
memorizing training patterns.

Taken together, on our benchmark GraphMCMNet behaves as a high-fidelity learned
approximation of the DP solver: 0.029\% error, zero validity violations, and a
measured ${>}100{\times}$ GPU-batched throughput advantage over sequential DP.
The validity gap is arguably the most consequential finding: regression-based cost
models~\cite{b17,b19} can silently output infeasible values in the majority of
cases, a risk that structural architectures eliminate by construction.

\section{Limitations and Threats to Validity}
\label{sec:limitations}

We state the boundaries of our claims explicitly.

\noindent\textbf{Data scope.} Training uses synthetic chains
($d{\in}[10,500]$). Although a preliminary evaluation on real architecture
chains (Section~\ref{sec:realarch}) suggests the conclusions transfer,
production workloads in compilers and query optimizers may exhibit dimension
statistics outside this range.

\noindent\textbf{Scale.} Evaluation covers $n{\le}50$. The GNN's sub-problem
graph has $O(n^2)$ nodes and $O(n^3)$ DP-child edges; memory becomes the
binding constraint for $n{\gg}100$.

\noindent\textbf{Single training run.} Reported metrics come from a single trained
checkpoint per model; bootstrap 95\% confidence intervals
(Table~\ref{tab:overall}) quantify evaluation variance but not training
variance across random seeds.

\noindent\textbf{Approximation, not replacement.} The learned models approximate
the DP solver; they do not replace it in settings requiring certified optimality.
The 93\% exact-match rate means 7\% of chains receive a slightly suboptimal
parenthesization.


%% ================================================================
%%   CONCLUSION AND REFERENCES
%% ================================================================

\section{Conclusion}
This paper studied the Matrix Chain Multiplication problem through the lens of
learned approximation, comparing structural reasoning---predicting the DP split
table and computing costs exactly---against statistical estimation, which regresses
costs directly. Evaluation across 500 diverse synthetic chains and a preliminary
real-architecture benchmark reveals a stark
dichotomy: regression-based methods fail systematically, whereas structural neural
models closely approximate the DP solver.

The most consequential observation is the \emph{mathematical validity gap}.
Despite 145 engineered features, tree-based regressors produce costs below the
theoretical DP minimum in roughly 80\% of cases, because they treat a
combinatorial objective as an unconstrained continuous variable. GraphMCMNet
instead predicts discrete split decisions, and---as we prove in
Theorem~\ref{thm:validity}---every such prediction corresponds to a feasible
parenthesization, yielding 100\% validity by construction. On our benchmark this
framing simultaneously reduces error to 0.029\% MAPE (95\% CI reported in
Table~\ref{tab:overall}).

Encoding DP sub-problem dependencies as graph edges also enables length
extrapolation: trained only on $n{\le}35$, GraphMCMNet achieves its lowest error
on chains of length 41--50. These findings support a broader thesis from neural
algorithmic reasoning~\cite{b25,b26,b27,b28}: embedding algorithmic structure
directly into neural architectures is a prerequisite for reliable learned
combinatorial optimization. Future work will scale the structural approach
to larger tensor computation graphs and evaluate on production compiler traces.

\begin{thebibliography}{00}

\bibitem{b1} S. S. Godbole, ``On efficient computation of matrix chain products,'' \emph{IEEE Trans. Computers}, vol. C-22, no. 9, pp. 864--866, Sep. 1973.
\bibitem{b2} T. C. Hu and M. T. Shing, ``Computation of matrix chain products. Part I \& II,'' \emph{SIAM J. Computing}, vol. 11, no. 2, pp. 362--373, 1982.
\bibitem{b3} F. L\'{o}pez, L. Karlsson, and P. Bientinesi, ``On the parenthesisations of matrix chains: All are useful, few are essential,'' \emph{arXiv:2303.17352}, 2023.
\bibitem{b4} T. H. Cormen, C. E. Leiserson, R. L. Rivest, and C. Stein, \emph{Introduction to Algorithms}, 3rd ed. Cambridge, MA: MIT Press, 2009.
\bibitem{b7} O. Vinyals, M. Fortunato, and N. Jaitly, ``Pointer Networks,'' in \emph{Proc. NeurIPS}, vol. 28, 2015.
\bibitem{b8} I. Bello, H. Pham, Q. V. Le, M. Norouzi, and S. Bengio, ``Neural combinatorial optimization with reinforcement learning,'' in \emph{Proc. ICLR}, 2017.
\bibitem{b9} D. Bahdanau, K. Cho, and Y. Bengio, ``Neural machine translation by jointly learning to align and translate,'' in \emph{Proc. ICLR}, 2015.
\bibitem{b11} A. Vaswani \emph{et al.}, ``Attention is all you need,'' in \emph{Proc. NeurIPS}, 2017.
\bibitem{b12} T. Chen and C. Guestrin, ``XGBoost: A scalable tree boosting system,'' in \emph{Proc. ACM SIGKDD}, pp. 785--794, 2016.
\bibitem{b13} L. Breiman, ``Random forests,'' \emph{Machine Learning}, vol. 45, pp. 5--32, 2001.
\bibitem{b17} R. Marcus \emph{et al.}, ``Neo: A learned query optimizer,'' \emph{Proc. VLDB Endowment}, vol. 12, no. 11, pp. 1705--1718, 2019.
\bibitem{b19} Z. Yang \emph{et al.}, ``Balsa: Learning a query optimizer without expert demonstrations,'' in \emph{Proc. ACM SIGMOD}, 2022.
\bibitem{b20} Q. Cappart \emph{et al.}, ``Combinatorial optimization and reasoning with graph neural networks,'' in \emph{Proc. IJCAI Survey Track}, 2021.
\bibitem{b21} Y. Bengio, A. Lodi, and A. Prouvost, ``Machine learning for combinatorial optimization: A methodological tour d'horizon,'' \emph{European J. Operational Research}, vol. 290, no. 2, pp. 405--421, 2021.
\bibitem{b22} N. Vesselinova, R. Steinert, D. F. Perez-Ramirez, and M. Boman, ``Learning combinatorial optimization on graphs,'' \emph{IEEE Access}, vol. 8, pp. 120388--120416, 2020.
\bibitem{b23} T. Chen \emph{et al.}, ``TVM: An automated end-to-end optimizing compiler for deep learning,'' in \emph{Proc. 13th USENIX OSDI}, 2018.
\bibitem{b24} R. Baghdadi \emph{et al.}, ``A deep learning based cost model for automatic code optimization,'' in \emph{Proc. MLSys}, 2021.
\bibitem{b5} A. Naz \emph{et al.}, ``Optimal sequence for chain matrix multiplication using evolutionary algorithm,'' \emph{PLOS ONE}, 2021.
\bibitem{b25} P. Veli\v{c}kovi\'{c} and C. Blundell, ``Neural algorithmic reasoning,'' \emph{Patterns}, vol.~2, no.~7, 2021.
\bibitem{b26} P. Veli\v{c}kovi\'{c} \emph{et al.}, ``The CLRS algorithmic reasoning benchmark,'' in \emph{Proc. ICML}, 2022.
\bibitem{b27} B. Ibarz \emph{et al.}, ``A generalist neural algorithmic learner,'' in \emph{Proc. Learning on Graphs (LoG)}, 2022.
\bibitem{b28} S. Mahdavi \emph{et al.}, ``Towards better out-of-distribution generalization of neural algorithmic reasoning tasks,'' \emph{Trans. Machine Learning Research (TMLR)}, 2023.
\bibitem{b29} B. Bevilacqua \emph{et al.}, ``Neural algorithmic reasoning with causal regularisation,'' in \emph{Proc. ICML}, 2024.

\end{thebibliography}

\end{document}
